% !TEX program = lualatex
% Transparencias de Fundamentos de las Matemáticas I
% Clase 05: Métodos de demostración

\documentclass[aspectratio=169,11pt]{beamer}

\usepackage{fontspec}
\usepackage[spanish,es-nodecimaldot]{babel}
\usepackage{amsmath,amssymb,mathtools}
\usepackage{booktabs}
\usepackage{csquotes}
\usepackage{xcolor}

\usetheme{Madrid}
\usecolortheme{default}
\setbeamertemplate{navigation symbols}{}
\setbeamertemplate{footline}[frame number]

\definecolor{FdMBlue}{RGB}{20,55,100}
\definecolor{FdMRed}{RGB}{130,30,30}
\definecolor{FdMGreen}{RGB}{25,100,60}

\setbeamercolor{structure}{fg=FdMBlue}
\setbeamercolor{title}{fg=white,bg=FdMBlue}
\setbeamercolor{frametitle}{fg=white,bg=FdMBlue}
\setbeamercolor{block title}{fg=white,bg=FdMBlue}
\setbeamercolor{block body}{bg=blue!3}
\setbeamercolor{alerted text}{fg=FdMRed}

\setmainfont{Latin Modern Roman}
\setsansfont{Latin Modern Sans}
\setmonofont{Latin Modern Mono}

\newcommand{\N}{\mathbb{N}}
\newcommand{\Z}{\mathbb{Z}}
\newcommand{\Q}{\mathbb{Q}}
\newcommand{\R}{\mathbb{R}}
\newcommand{\C}{\mathbb{C}}
\newcommand{\Pow}{\mathcal{P}}
\newcommand{\id}{\operatorname{id}}
\newcommand{\mcd}{\operatorname{mcd}}
\newcommand{\mcm}{\operatorname{mcm}}

\title[FdM I -- Clase 05]{Métodos de demostración}
\subtitle{Directa, contraposición, contradicción y casos}
\author{Antonio Falcó Montesinos}
\institute{Universidad CEU Cardenal Herrera}
\date{Fundamentos de las Matemáticas I}

\begin{document}

\begin{frame}
  \titlepage
\end{frame}

\begin{frame}{Objetivos de la clase}
\begin{itemize}
  \item Elegir un método de demostración según la forma del enunciado.
  \item Redactar pruebas directas y por contraposición.
  \item Reconocer cuándo usar contradicción o división en casos.
\end{itemize}
\end{frame}

\begin{frame}{Mapa de la clase}
\tableofcontents
\end{frame}

\section{Demostración directa}

\begin{frame}{Demostración directa}
\begin{block}{Idea}
\begin{itemize}
  \item Se parte de las hipótesis y se llega a la conclusión mediante pasos justificados.
  \item Es el método natural para muchas implicaciones.
  \item Cada paso debe apoyarse en definiciones o resultados previos.
\end{itemize}
\end{block}
\end{frame}

\begin{frame}{Demostración directa}
\begin{block}{Modelo}
\begin{itemize}
  \item Para demostrar \(P\Rightarrow Q\): supongamos \(P\).
  \item Usando definiciones, transformamos la información de \(P\).
  \item Concluimos \(Q\).
\end{itemize}
\end{block}
\end{frame}

\section{Métodos indirectos}

\begin{frame}{Métodos indirectos}
\begin{block}{Contraposición}
\begin{itemize}
  \item Para probar \(P\Rightarrow Q\), se prueba \(\neg Q\Rightarrow\neg P\).
  \item Es útil cuando la negación de la conclusión aporta información concreta.
  \item Ejemplo típico: si \(n^2\) es par, entonces \(n\) es par.
\end{itemize}
\end{block}
\end{frame}

\begin{frame}{Métodos indirectos}
\begin{block}{Contradicción}
\begin{itemize}
  \item Se supone lo contrario de lo que se quiere demostrar.
  \item A partir de esa suposición se obtiene una contradicción explícita.
  \item No basta decir ``supongamos lo contrario'': hay que llegar a una imposibilidad.
\end{itemize}
\end{block}
\end{frame}

\begin{frame}{Métodos indirectos}
\begin{block}{Demostración por casos}
\begin{itemize}
  \item Se divide el problema en situaciones exhaustivas y excluyentes o manejables.
  \item Cada caso debe cubrirse.
  \item Al final se concluye que el resultado vale en todos los casos.
\end{itemize}
\end{block}
\end{frame}

\section{Equivalencias}

\begin{frame}{Equivalencias}
\begin{block}{Si y solo si}
\begin{itemize}
  \item Para demostrar \(P\Leftrightarrow Q\), se prueban dos implicaciones.
  \item Primero \(P\Rightarrow Q\).
  \item Después \(Q\Rightarrow P\).
\end{itemize}
\end{block}
\end{frame}

\section{Profundización}

\begin{frame}{Profundización}
\begin{block}{Elegir método}
\begin{itemize}
  \item Si el enunciado es universal, se empieza con un elemento arbitrario.
  \item Si el enunciado es existencial, se construye un objeto concreto.
  \item Si la conclusión negada es más manejable, puede convenir la contraposición.
\end{itemize}
\end{block}
\end{frame}

\begin{frame}{Profundización}
\begin{block}{Modelo por contraposición}
\begin{itemize}
  \item Objetivo: si \(n^2\) es par, entonces \(n\) es par.
  \item Contrapositiva: si \(n\) no es par, entonces \(n^2\) no es par.
  \item La prueba usa la forma \(n=2k+1\) para enteros impares.
\end{itemize}
\end{block}
\end{frame}

\begin{frame}{Profundización}
\begin{block}{Errores de redacción}
\begin{itemize}
  \item Usar ejemplos como si fueran una prueba general.
  \item No indicar dónde se usa una hipótesis.
  \item Terminar con una igualdad correcta pero sin conectar con la conclusión.
\end{itemize}
\end{block}
\end{frame}

\begin{frame}{Cierre}
\begin{itemize}
  \item La forma del enunciado sugiere el método.
  \item Una prueba bien redactada declara hipótesis y objetivo.
  \item Los ejemplos orientan, pero la demostración cierra.
\end{itemize}
\end{frame}

\begin{frame}{Trabajo recomendado}
\begin{itemize}
  \item Releer las secciones correspondientes del manual.
  \item Identificar definiciones, símbolos nuevos y enunciados principales.
  \item Preparar dudas y ejemplos para el seminario asociado.
\end{itemize}
\end{frame}

\end{document}
